\color{uniblue}{\section{Утилизация}}
\color{black}

После  окончания  срока  службы  устройство  подлежит  демонтажу  и  утилизации  согласно ГОСТ~Р~55102. Изделие не содержит в своём составе опасных или ядовитых веществ, способных нанести вред здоровью человека или окружающей среде и не представляет опасности для жизни, здоровья людей 
и окружающей среды по окончании срока службы. В этой связи утилизация изделия может производиться по правилам утилизации общепромышленных отходов. 
Содержание драгоценных металлов в компонентах изделия (электронных платах, разъёмах и т.п.) крайне мало, поэтому их вторичную переработку производить нецелесообразно. Демонтаж  и  утилизация  устройства  не  требуют  применения  специальных  мер  безопасности  и выполняются без применения специальных приспособлений и инструментов. Утилизацию должны производить следующие организации:
\begin{itemize} 
\item производители электротехнического и электронного оборудования; 
\item предприятия по переработке ОЭЭО; 
\item специализированные пункты сбора и хранения ОЭЭО; 
\item пункты сбора вторичного сырья.
\end{itemize} 
Основным  методом  утилизации  должна  быть  разборка  устройства  по  группам  и  обращению  с отходами согласно ГОСТ~Р~55102. 